\scanpage{11}
\[
G''\supset H''\quad{}^{Q}H'\subset G'\quad{}^{P}G_k
\qquad
\boxed{\begin{aligned}
\alpha&\in P(\mathbb C)\\
\beta&\in Q(\mathbb C)\subset P'(\mathbb C)\\
\gamma&\in \deleted{\unsure{R}}P''(\mathbb C)
\end{aligned}}
\qquad
\begin{aligned}
\beta\alpha&=\gamma^{-1}\\
\gamma\beta&=\alpha^{-1}\\
\alpha\gamma&=\beta^{-1}
\end{aligned}
\]
\[
\left.\begin{aligned}
P&=\underline{\mathrm{Isom}}(T_B\otimes k,T_{\mathrm{DR}})\\
Q\subset P'&=\underline{\mathrm{Isom}}(T_{\mathrm{DR}},T_{\mathrm{Hdg}})\simeq Q\mathbin{\overset{H'}{\times}}G'\\
P''&=\underline{\mathrm{Isom}}(T_{\mathrm{Hdg}},T_B\otimes k)
\end{aligned}\right\}\text{ déf}/k
\]
\diagram{\begin{tikzcd}[ampersand replacement=\&,column sep=small]
\& T_{\mathrm{Hdg}}\otimes_k\mathbb C\arrow[dr,"\gamma"] \&\\
T_{\mathrm{DR}}\otimes_k\mathbb C\arrow[ur,"\beta"] \& \& T_B\otimes_{Q}\mathbb C\arrow[ll,"\alpha"]
\end{tikzcd}}

\[
\underbrace{U'}_{\substack{\text{rad.}\\\text{unipotent}\\\text{de }\unsure{U'|B'}}}
\subset\underbrace{B'}_{\substack{\text{parabolique}\\\text{associé à la}\\\text{filtration de}\\T_{\mathrm{DR}}}}\subset G'
\]
$Q=$ fibré des sous groupes de Lévi de $B'$
\[
 B'/U'\simeq L'\xleftarrow[\substack{\text{hom}\\\text{central}\\(\mathrm{Tate})}]{j'}G_m
\]

\[
\left.\begin{aligned}
\mathrm{Fil}&=\underline{\mathrm{Filt}}_{\substack{\otimes,\,\text{dans }\deleted{\unsure{\mathcal F}}\\
\text{classe int. de}\\\text{la filtration}\\\text{de }T_{\mathrm{DR}}}}\;T_{\unsure{b|B}}\\
\mathrm{Big}&=\underline{\mathrm{Bigr}}_{\substack{\otimes,\,\text{dans}\\\text{classe int de}\\
\text{la }\deleted{\unsure{\text{filtr.}}}\text{bigraduation}\\\text{de }T_{\mathrm{DR}}}}\;T_B
\end{aligned}\right\}\text{ déf}/Q
\]
\margin{$\nwarrow$ NB Soit $F\supset\mathrm{Fil},\Phi\supset\mathrm{Big}$ définis comme ci-contre, mais sans condition de forme intérieure. Alors $G$ y opère, et $A=F/G,B=\Phi/G$ sont des schémas \uline{\unsure{indét.|étales}} sur $Q$. Alors $\varphi\in F(\mathbb C),\psi\in\Phi(\mathbb C)$ \deleted{\unsure{définissent}} des pts $\dot\varphi\in A(\mathbb C),\dot\psi\in B(\mathbb C)$. \deleted{\unsure{$k$}} Soient $\varphi_1\in A,\psi_1\in B$ leurs images dans $A$ et $B$ (\unsure{pts} fermés \unsure{induits|susdits}) Par définition, $\mathrm{Fil}=F_{\varphi_1},\mathrm{Big}=\Phi_{\psi_1}$}
\margin{\uline{NB} Si on a $\varepsilon:\deleted{\unsure{\mathrm{Soit}}}G\longrightarrow G_m$ on a $(\varepsilon|(B'/U'))\circ j=\mathrm{id}.$ donc $j$ \unsure{mon.} Si $H'\subset B'$ est l'image inverse de $\operatorname{Im}j'$ dans $B'$, \deleted{\unsure{le tore}} $j'$ est \inserted{donc} connu quand on connait $H'$}

\[
P\times P'\longrightarrow P'',\quad P'\times P''\unsure{\mathbin{-}|\longrightarrow}P,\quad
P''\times P\longrightarrow P',\qquad Q\subset P'
\]

$\mathrm{Big}\xrightarrow[\deleted{\unsure{d}}]{p}\mathrm{Fil}$ fibration, faisant de $B$ un torseur sur $\mathrm{Fil}$, de groupe $U_{\mathrm{Fil}}$, un groupe unipotent défini sur $F$, de dim. relative égale à $\dim U'$. $\mathrm{Fil}$ est \unsure{des|de} dim égale à $\dim G'/B'$, $\mathrm{Big}$ est de dim égale à $\dim G'/B'+\dim U'=\dim G'/L'$, où $L'\deleted{\unsure{=}}=\operatorname{Centr}(j'_0)$, $j'_0:G_m\xrightarrow[k]{}H'$ relevant $j'$.

\margin{en fait, $\dot\varphi,\dot\psi$ sont dans $A(k),B(k)$ donc dans $A(k_1),B(k_1)$, où $k_1$ est la clôture alg. de $Q$ dans $k$. D'ailleurs, \deleted{$F\longrightarrow\Phi$ \unsure{induit}} $A\xrightarrow{\sim} B$, donc}
\margin{\unsure{constr.} explicite en termes du groupe déployé, associé au réseau \unsure{cor.} $\Gamma$ avec $W$ (Weyl) agissant : on \unsure{tord|aura} $\Gamma/W$ par autom. extérieurs\ldots.}

On a un diagramme commutatif
\diagram{\begin{tikzcd}[ampersand replacement=\&]
P''\arrow[d,"{\deleted{\unsure{q}}r}"'] \& P\arrow[d,"{q\deleted{\unsure{*}}}"]\\
\mathrm{Big}_k\arrow[r,"p_k"'] \& \mathrm{Fil}_k
\end{tikzcd}}
\[
\boxed{\begin{aligned}
\varphi&=q(\alpha)\in\mathrm{Fil}(\mathbb C)\\
\psi&=r(\gamma)\in\mathrm{Big}(\mathbb C)
\end{aligned}}
\qquad [\varphi=p(\psi)]
\qquad\begin{aligned}
\dim G&=g\\\dim U'&=u.\\\dim B'&=b\\\dim L'&=l\\
(u+l&=b)
\end{aligned}
\]

\diagram{\resizebox{\linewidth}{!}{\fbox{\begin{tikzcd}[ampersand replacement=\&,column sep=small,row sep=small]
\& {k(\alpha,\beta)=k(\beta,\gamma)=k(\gamma,\alpha)=K}
 \arrow[dl,dash,"{\leq u}"']\arrow[d,dash,"{\leq u}"]\arrow[dr,dash,"{\leq g}"] \& \&\\
k(\gamma)\arrow[d,dash,"{\leq l}"] \& k(\alpha)\arrow[d,dash,"{\leq b}"] \& k(\beta)\arrow[d,dash,"{\leq u}"] \&\\
k(\psi)\arrow[r,dash,"{\leq u}"]\arrow[d,dash,"{\text{comp}}"] \& k(\varphi)\arrow[r,dash,"{\leq g-b}"]\arrow[d,dash,"{\text{comp}}"] \& k\arrow[d,dash]\arrow[dr,dash] \&\\
Q(\psi)\arrow[r,dash,"{\leq u}"'] \& Q(\varphi)\arrow[r,dash,"{\leq g-b}"']\arrow[d,dash] \& Q \arrow[r,dash,"{\mathrm{alg}}"'] \& \deleted{\unsure{Q(\psi_1)}}\,\inserted{Q(\varphi_1)}\unsure{=|\simeq}Q(\varphi_1)\\
\&\deleted{\unsure{Q}} \& \&
\end{tikzcd}}}}
\margin{$Q(\varphi_1)$ est \uline{indépendant} du plongement de $k$ dans $\mathbb C$, donc \unsure{ce corps} doit être \unsure{le corps de déf.|un inv. du gr.} du motif}
\[
\deleted{\unsure{k(\alpha,\beta)}}K\subset\deleted{\unsure{K}}Q(k(\alpha),\overline{k(\alpha)})
\]
$(2i\pi)^n\in k(\alpha)$ si $Q(-n)\in\operatorname{Ob}\mathcal M$

\uline{Conj. de transcendance forte}

Si $k\subset\overline{\unsure{Q|\mathbb Q}}$, alors \unsure{?} $\operatorname{degr.tr.}\underbrace{k(\alpha,\beta)}_{K}/k=g+u$
\[
\Longrightarrow\left\{\begin{aligned}
\operatorname{deg.tr.} k(\alpha)/\unsure{k|K}&=g\\
\operatorname{deg.tr.} k(\beta)/\unsure{K|k}&=u\\
\operatorname{degr.} k(\gamma)/k&=g\\
\operatorname{deg.tr.} K/k(\beta)&=g,\quad\text{et}\\
\operatorname{deg.tr.} K/k(\alpha)&=\operatorname{deg\,tr}K/k(\gamma)=u\\
\operatorname{deg\,tr} Q(\varphi)/Q&=g-b\\
\operatorname{deg\,tr.} Q(\psi)/Q&=g-b+u=g-l
\end{aligned}\right.
\]
% archive sheet number: 10
